\section{Method}\label{sec:method}
% Written in P8 (framework-independent). Describes only what the code does: src/data/{features,windows,splits}.py,
% src/models/*.py, src/calibration.py, src/quantize.py, configs/{data,models,quantization}.yaml, quant_meta.json;
% decisions K-002, K-007, K-016 ... K-025, K-027 ... K-045. No numeric literal: every number is a macro.

Fig.~\ref{fig:system} summarises the pipeline.

\begin{figure*}[t]
\centering
\includegraphics[width=\textwidth]{fig1_system}
\caption{System overview. CAN captures are split by source capture (ROAD) or by disjoint time blocks
(Car-Hacking), encoded per frame and cut into \WinSize-frame windows (non-overlapping at validation and test
time). A 1D-CNN, a convolutional autoencoder or XGBoost scores each window; the two deep models are post-training
quantized to fp16 (\Qhalf) and INT8 (\Qint), each version is thresholded at a \TargetFPR{} false-positive rate on
attack-free validation windows, and every version is evaluated per attack type.}
\label{fig:system}
\end{figure*}

\subsection{Frame encoding}
Each frame becomes a \NFeatures-dimensional vector in the unit interval:
\begin{itemize}
\item an \emph{identifier code}: the rank of the identifier in the sorted vocabulary of identifiers seen in benign
training frames, divided by the vocabulary size $K$ ($K=\VocabCH$ for Car-Hacking, $K=\VocabROAD$ for ROAD);
identifiers outside the vocabulary, such as most injected fuzzing identifiers, map to a reserved code zero;
\item the \NPayloadBytes{} \emph{payload bytes}, each divided by \ByteScale; shorter payloads are zero-padded;
\item the \emph{inter-arrival time} since the previous frame with the same identifier, on a logarithmic scale
between a floor and a cap, the \IATFloorQ{} and \IATCapQ{} quantiles of the benign per-identifier inter-arrival
times of the training segments (Car-Hacking \IATFloorUsCH~$\mu$s to \IATCapSCH~s; ROAD \IATFloorUsROAD~$\mu$s to
\IATCapSROAD~s). It is computed inside each segment, and the first frame of an identifier in a segment takes the
cap value.
\end{itemize}
The vocabulary, floor and cap are fitted on training data only. The encoding is designed for the INT8 input: the
asymmetric unsigned INT8 input quantizer of every \Qint{} graph has scale \InputScale{} and zero point
\InputZeroPoint, that is \InputLevels{} levels on the unit interval. The identifier codes, including the reserved
code (\IdLevelsCH{} and \IdLevelsROAD{} values), and the payload-byte values fall on distinct levels, which we check for every quantized graph, whereas scaling raw
identifiers to the unit interval would merge neighbouring identifiers into one level. The logarithmic
inter-arrival time gives every decade the same resolution; a linear scale would squeeze the typical CAN periods
into a few levels. A loss measured after quantization therefore reflects the quantized model, not the input
scaling.

\subsection{Windows and labels}
A window holds $N=\WinSize$ consecutive frames and is labelled \emph{attack} if it contains at least one injected
frame. Windows are cut with stride \StrideTrain{} in training (\OverlapTrainPct{} overlap, as data augmentation)
and with stride \StrideEval, without overlap, in validation and test.

\subsection{Leakage-safe split}
Random splitting of windows overestimates performance on CAN data for several reasons, and the split is designed
against each of them.
% TODO-CITE: overlapping windows and random window-level splits in CAN IDS evaluation.
\begin{itemize}
\item \emph{Overlapping windows.} Two consecutive training windows share \OverlapTrainPct{} of their frames. If
windows of one recording were assigned to splits at random, almost every test window would have a near copy in
training. We therefore split recordings (or time blocks of a recording) before windowing, and overlap exists only
inside the training split.
\item \emph{Non-overlapping test windows.} With stride equal to the window size, every test frame is evaluated
exactly once, so an injection event is not counted several times and test windows do not share frames.
\item \emph{Duplicated recordings.} Two kinds of recordings occur twice in the data. ROAD masquerade captures are
copies of their fabrication source (Section~\ref{sec:threat}); the unit of the ROAD split is therefore the source
capture, and both versions always land in the same split. The \SharedTailFiles{} Car-Hacking attack files end with
the same attack-free recording (identical timestamps, identifiers and payloads); it is cut from every file and used
once, in training.
\item \emph{Fitted parameters.} The identifier vocabulary and the inter-arrival-time scaling are fitted on training
data, decision thresholds on validation data and INT8 calibration on training data; test data are never used for
fitting or selection.
\end{itemize}

\emph{ROAD} has several captures per attack type. For every fabrication type the first, second and third source
capture go to training, validation and test, with the derived masquerade captures. Two exceptions follow from
window counts alone (at least \MinTestAttackWindows{} test attack windows per type): the first fuzzing capture is
used for training and the other two for testing, without a validation capture, and the engine-coolant pair, which
has a single source capture, is placed in the test split as an attack type unseen in training and reported apart.
Attack-free captures are distributed over the three splits; \ExcludedROAD{} captures without frame-level labels
are excluded.

\emph{Car-Hacking} has one file per attack type, so a file-level split would remove each type from training or from
test. Each attack file is cut in time instead: the first \TrainBursts{} attack bursts go to training, the next
\ValBursts{} to validation and the remaining bursts to test, and \GapBursts{} whole bursts are discarded at each
cut. The attack-free file is split by time (\NormalFractions{} percent) with \NormalGapS-s gaps. Every gap is at
least \MinBlockGapS~s long, an order of magnitude longer than the inter-arrival-time cap, so no feature of one block
looks into the next. Blocks of one file are still parts of the same driving session, which is a weaker
independence than a file-level split (Section~\ref{sec:limitations}).

An automated test enforces the split: no sequence of \LeakNgram{} consecutive frames (identifiers, payloads and
inter-frame time differences, hence independent of timestamp offsets) occurs in two splits, and it detects the
known duplicates (masquerade--fabrication pairs, the shared Car-Hacking recording) as positive controls. The same
test checks that the test stride equals the window size and that the feature scaling was fitted on training data
only.

\subsection{Models}
\emph{1D-CNN} (supervised). The network stacks \CNNBlocks{} convolution blocks (channels \CNNChannels; kernel sizes
\CNNKernels), each with batch normalisation and ReLU and with max pooling between blocks, followed by an average over time and a
two-layer head with \CNNHidden{} hidden units and dropout \CNNDropout. The window score is the output logit
(\CNNParams{} parameters).

\emph{Convolutional autoencoder} (unsupervised, conv-AE). Two strided convolutions (channels \AEChannels) and a
convolution to \AELatentChannels{} channels encode a window to a latent sequence of length \AELatentLen; two
transposed convolutions and a convolution with a sigmoid output reconstruct it (\AEParams{} parameters). It is
trained on attack-free windows only, and the window score is the mean squared reconstruction error.

The planned unsupervised model was an LSTM autoencoder (\LSTMHidden{} hidden units, latent size \LSTMLatent). A
feasibility rule fixed before training required a validation ROC-AUC of at least \FeasMinAUC{} and allowed, on
failure, one retry with \RetryHidden{} hidden units, learning rate \RetryLR{} and up to \RetryMaxEpochs{} epochs,
followed by the conv-AE. The candidate order and the selection rule (the first candidate that passes on both
datasets is used) were fixed before the retry and the conv-AE were evaluated. The LSTM autoencoder failed on ROAD (Car-Hacking \FeasLSTMBaseCH, ROAD \FeasLSTMBaseROAD; retry
\FeasLSTMRetryCH{} and \FeasLSTMRetryROAD), and the conv-AE passed (\FeasConvAECH{} and \FeasConvAEROAD). Validation
attack labels were used for this pass/fail decision only, never for a threshold.

\emph{XGBoost} (supervised reference). Gradient-boosted trees operate on \NSummaryFeatures{} window summaries
computed from the same encoded windows: identifier entropy, number of distinct identifiers and share of
out-of-vocabulary identifiers; mean, standard deviation, minimum and lower decile of the inter-arrival time; mean
and variance of every payload byte; and the share of consecutive same-identifier frames whose payload changes. The
score is the raw margin. XGBoost is not quantized; it serves as a reference for size, latency and detection.

\subsection{Decision threshold}
All models use the same rule. The threshold is the \CalibQuantile{} quantile of the scores of the \emph{clean
validation normals}: validation windows labelled normal that do not come from a masquerade capture and have no
frame within \NormalMarginS~s of an injected frame. A window is an attack if its score exceeds the threshold, so
the false-positive rate on these windows is at most \TargetFPR. No attack label and no test window is used. Fixing
the false-alarm rate compares all models and versions at the same operating point, which corresponds to the
false-alarm budget of an in-vehicle detector. Every quantized version receives its own threshold from its own
validation scores (primary analysis). As a secondary analysis, the float32 threshold is applied unchanged to the
quantized version to measure the \emph{calibration shift} that a deployment without re-calibration would see.

\subsection{Quantization}
\emph{\Qfp{} (float32)} is the PyTorch model exported to ONNX (opset \Opset), checked against PyTorch on the CPU
(absolute tolerance \QZeroAtol, relative tolerance \QZeroRtol). All quantization losses are computed relative to
this ONNX model on the CPU, so that the reference and the quantized versions run in the same runtime.

\emph{\Qhalf{} (float16)} converts the ONNX graph to half precision with float32 inputs and outputs. The CPU
execution provider of ONNX Runtime has no float16 kernel for the heavy layers of our models (\HeavyFPCNN{} of
\HeavyCNN{} for the 1D-CNN, \HeavyFPAE{} of \HeavyAE{} for the conv-AE) and casts them back to float32. A rule set
before any float16 result was seen states that the CPU file does not represent float16 if less than
\MinHeavyFPShare{} of the heavy layers run in half precision. The float16 accuracy is therefore measured with the
PyTorch model in half precision (weights, inputs and activations) on the GPU, the model size is that of the float16
ONNX file, and no CPU latency is reported for \Qhalf.

\emph{\Qint{} (INT8)} is static post-training quantization with ONNX Runtime in the QDQ format: asymmetric unsigned INT8
activations, symmetric per-channel signed INT8 weights and min--max calibration on \CalibWindowsQ{} windows drawn with
the run seed from the model's own training selection (1D-CNN: supervised training windows; conv-AE: attack-free
training windows). The code refuses calibration windows from outside the training split or fewer than
\MinCalibWindowsQ. The quantized operator types are \QOpTypes; the remaining nodes (transpositions, reductions and
the score computation) stay in float, and the list of float nodes is stored with every model. A kernel profile
records how each heavy layer runs: all \HeavyIntCNN{} heavy layers of the 1D-CNN use INT8 kernels, while for the
conv-AE \HeavyIntAE{} of \HeavyAE{} do; its \HeavySimAE{} transposed convolutions have no INT8 CPU kernel and run as
dequantize--float--quantize simulation, which keeps the INT8 numerics but not the INT8 arithmetic. The score (logit
or reconstruction error) is computed inside every graph, so latency covers the whole inference.
